% !TeX root = ../../Book5.tex
% 纲要标题：### 9.5 基本代数、Ext 箭图与关系
% 纲要位置：工作记录/计划纲要.md，第 4366 行。

\section{基本代数、Ext 箭图与关系}
\label{b5:sec:0905}


% 纲要范围：
%
% * 对代数闭域上的有限维代数，调用第二卷的投射覆盖和 Morita 理论构造基本代数，并说明分裂基域假设的用途
% * 从本原幂等元及 \(J/J^2\) 构造箭图，按第9.2节既定的路径和左模约定核验箭的方向及 \(\operatorname{Ext}^1\) 维数
% * 证明基本代数可写成路径代数模容许关系理想，并用箭理想的幂界解释容许性和有限维性
% * 以三角矩阵代数、双数代数及一个带关系例子连接第二卷的根滤过；一般有关系代数不直接适用第10章无关系箭图的分类
% * 从第9章可直接进入第10章；所需 ADE 二次型和反射将在第10.1节就地建立，不要求先读 Lie 代数章节
%
% ---
%

前面从箭图造出代数。现在反过来：一个有限维代数的简单模之间有哪些最初的联系，能否由箭表示？答案要去掉半单商中重复的矩阵大小，以根的第一层作箭，以更高层的乘法作关系。利用第二卷的根、幂等元提升与 Morita 等价，本节将代数闭域上的有限维代数化为与它 Morita 等价的带关系路径代数 \(kQ/I\)。\chapref{b5:ch:10}研究的则是有限无环箭图的无关系路径代数 \(kQ\)，其有限维表示具有\cref{b5:thm:acyclic-quiver-resolution}给出的长度一投射消解；加入关系以后，这种消解及不可分解表示的分类都可能改变。

\subsection{基本代数与分裂基域假设}
\label{b5:subsec:0905:01}

% 跨卷引用目标：b2:thm:finite-algebra-radical-nilpotent（第二卷定理15.1.5）。
% 跨卷引用目标：b2:thm:orthogonal-idempotent-lifting（第二卷定理15.3.4）。
% 跨卷引用目标：b2:thm:artin-wedderburn（第二卷定理13.3.1）。
设 \(A\) 为代数闭域 \(k\) 上的有限维含幺代数，记 \(J=J(A)\)。第二卷定理15.1.5证明了 \(J\) 幂零；第二卷定理13.3.1的 Artin--Wedderburn 结构定理把商写为
\[
A/J\cong\prod_{i=1}^r M_{n_i}(k).
\]
每个简单模的自同态除环在代数闭域上为 \(k\)，所以这里没有额外的除代数系数。第二卷定理15.3.4又允许将各矩阵块的对角矩阵单位提升为一组完整正交幂等元。
\begin{definition}\label{b5:def:split-basic-algebra}
设 \(k\) 为域，\(B\) 为非零有限维含幺 \(k\) 代数。若
\(B/J(B)\cong k^r\) 对某个 \(r\ge1\) 成立，就称 \(B\) 为一个\term[fenlie-jiben-daishu]{分裂基本代数}[split basic algebra]。
\end{definition}
第二卷第15.4节定义的一般基本代数只要求半单商为有限个除 \(k\) 代数的直积。在代数闭域上，这些有限维除代数均为 \(k\)，所以每个基本代数都是分裂基本代数。
\begin{proposition}\label{b5:prop:basic-morita-reduction}
设 \(k\) 为代数闭域，\(A\) 为非零有限维含幺 \(k\) 代数。可取一个幂等元 \(e\in A\)，使 \(AeA=A\)，且 \(B=eAe\) 是分裂基本代数。将 \(Ae\) 按乘法看成 \((A,B)\)-双模，则函子
\[
M\longmapsto eM,\qquad N\longmapsto Ae\otimes_B N
\]
给出左 \(A\) 模与左 \(B\) 模的等价，并保持有限维性。
\end{proposition}
\begin{proof}
% 跨卷引用目标：b2:prop:basic-corner（第二卷命题15.4.5）。
在每个 \(M_{n_i}(k)\) 块中选一个秩一对角幂等元，取相容提升 \(e_i\)，并令 \(e=\sum_i e_i\)。第二卷命题15.4.5关于基本角代数的结论在这里给出
\[
J(eAe)=eJe,\qquad eAe/eJe\cong k^r.
\]
这些条件也可直接检查：\(eJe\) 幂零，商角环恰是在每个矩阵块保留一个对角位置，因此半单且为 \(k^r\)。
% 跨卷引用目标：b2:prop:full-idempotent-morita（第二卷命题16.3.6）。
商代数中每个所选秩一幂等元生成整个对应矩阵块，故 \(AeA+J=A\)。令 \(M=A/AeA\)，则 \(JM=M\)；迭代得到 \(M=J^NM=0\)，所以 \(AeA=A\)。这验证了第二卷命题16.3.6的满幂等元条件。对左 \(A\)-模 \(M\)，该命题证明乘法映射
\[
Ae\otimes_B eM\longrightarrow M,\qquad ae\otimes m\longmapsto am
\]
为同构；对左 \(B\)-模 \(N\)，另一复合函子的值则经
\[
e(Ae\otimes_BN)\cong eAe\otimes_BN
=B\otimes_BN\xrightarrow{\ \sim\ }N
\]
同构于 \(N\)。第一个识别由右 \(B\)-模的分裂投影 \(Ae\to eAe\)、\(p\mapsto ep\) 诱导，最后一个映射为 \(b\otimes n\mapsto bn\)。这些同构对模同态相容，证明所列函子互为拟逆。这里 \(Ae\) 有限维，故两个函子也保持有限维对象。
\end{proof}
\begin{example}\label{b5:exm:matrix-basic-reduction}
设 \(k\) 为域，\(A=M_2(k)\)、\(e=E_{11}\)，其中 \(E_{ij}\) 为矩阵单位。则 \(eAe=kE_{11}\cong k\)，\(AeA=A\)。对有限维左 \(A\)-模 \(M\)，基本约化取出向量空间 \(eM\)。在 \(Ae\otimes_k eM\) 的第一个因子上取左 \(A\) 作用，则恢复映射与逆为
\[
\begin{aligned}
\mu_M:Ae\otimes_k eM&\longrightarrow M,
&p\otimes w&\longmapsto pw,\\
\delta_M:M&\longrightarrow Ae\otimes_k eM,
&m&\longmapsto\sum_{i=1}^2E_{i1}\otimes E_{1i}m.
\end{aligned}
\]
它们是相互逆的自然左 \(A\)-模同构。因此有限维 \(k\) 向量空间 \(W\) 对应于 \(Ae\otimes_kW\cong k^2\otimes_kW\)，并且
\[
\dim_k M=2\dim_k eM.
\]
\end{example}
\begin{proof}
对全部 \(i,j\in\{1,2\}\)，有 \(E_{ij}=E_{i1}eE_{1j}\)，故 \(AeA\) 包含全部矩阵单位。这说明 \(e\) 满，虽然它不等于 \(1_A\)。右 \(eAe\cong k\) 模 \(Ae\) 以 \(E_{11},E_{21}\) 为基；每个 \(E_{1i}m\) 都属于 \(eM\)，所以逆公式的两项确实落在所写张量积中。
求值映射 \(\mu_M\) 左 \(A\) 线性，而
\[
\mu_M\delta_M(m)=\sum_{i=1}^2 E_{i1}E_{1i}m
=\sum_{i=1}^2E_{ii}m=m.
\]
反向复合只需在 \(E_{i1}\otimes w\)、\(w\in eM\) 上计算：
\[
\delta_M\mu_M(E_{i1}\otimes w)
=\sum_{j=1}^2E_{j1}\otimes E_{1j}E_{i1}w
=E_{i1}\otimes E_{11}w=E_{i1}\otimes w.
\]
因此 \(\delta_M\) 是 \(\mu_M\) 的逆，也左 \(A\) 线性。模同态与矩阵单位的作用交换，所以两个公式都与模同态相容。对向量空间 \(W\)，另一复合给出 \(e(Ae\otimes_kW)=kE_{11}\otimes_kW\cong W\)，而 \(Ae\) 的二维基给出所列维数公式。
\end{proof}
例如自然简单模 \(k^2\) 被送到一维的第一坐标轴 \(e(k^2)\)，再由 \(Ae\otimes_k e(k^2)\) 恢复。Morita 等价保持模同态、直和和短正合列，从而保持简单性、不可分解性及扩张是否分裂；它不保持底层 \(k\) 维数。在一般证明中，每个矩阵块只保留一个秩一幂等元，正是因为同一块的列模同构，它们的投射覆盖也同构，只需保留每种类型的一份。

\ProseParagraphBreak

非代数闭域上的限制不能忽略。例如实代数 \(\mathbb C\) 已是除代数，只有一个简单模类型；一个普通单顶点、无箭的实箭图只给出 \(\mathbb R\)，不能用它记录 \(\mathbb C\) 的除代数结构。本节的普通箭图结论以分裂条件为基础。

\subsection{根商构造的箭图、箭向与扩张维数}
\label{b5:subsec:0905:02}

以下设 \(B/J\cong k^r\)，并选取提升标准坐标元的一组完整正交幂等元 \(e_1,\ldots,e_r\)，其中 \(J=J(B)\)。简单模为 \(S_i=Be_i/Je_i\)，每个都一维。选取每个空间
\[
e_j(J/J^2)e_i
\]
的一组 \(k\) 基，对每个基元画一支箭 \(i\to j\)。这里包括 \(i=j\) 的圈。根只保留第一层，是因为 \(J^2\) 中的元素已经可能由两支箭的复合产生，不应当再次作为独立箭。
\begin{proposition}\label{b5:prop:ext-quiver-direction}
设 \(k\) 为域，\(B\) 为非零有限维含幺分裂基本 \(k\) 代数，\(J=J(B)\)，\(B/J\cong k^r\)。取一组完整正交幂等元 \(e_1,\ldots,e_r\in B\)，其模 \(J\) 的像分别为 \(k^r\) 的标准坐标幂等元，令简单左模 \(S_i=Be_i/Je_i\)。则对 \(1\le i,j\le r\) 有
\[
\begin{aligned}
\operatorname{Ext}_B^1(S_i,S_j)
&\cong \operatorname{Hom}_{B/J}\bigl(Je_i/J^2e_i,S_j\bigr)\\
&\cong \operatorname{Hom}_k\bigl(e_j(J/J^2)e_i,k\bigr).
\end{aligned}
\]
后一同构使用 \(S_j=Be_j/Je_j\) 的指定向量 \(e_j+Je_j\)，以 \(c(e_j+Je_j)\mapsto c\) 识别 \(S_j\cong k\)。
以 \(1,\ldots,r\) 为顶点，对每个 \(e_j(J/J^2)e_i\) 的基元画一支箭 \(i\to j\)，则所得箭图中从 \(i\) 到 \(j\) 的箭数为 \(\dim_k\operatorname{Ext}_B^1(S_i,S_j)\)。
\end{proposition}
\begin{proof}
\(Be_i\) 是正则模的直和因子，因而投射。将\cref{b5:prop:projective-presentation-extensions}用于
\(0\to Je_i\to Be_i\to S_i\to0\)。
从 \(Be_i\) 到 \(S_j\) 的任何模同态都把 \(Je_i\) 送到 \(JS_j=0\)，故该 Hom 映射的限制为零。于是
\[
\operatorname{Ext}_B^1(S_i,S_j)\cong\operatorname{Hom}_B(Je_i,S_j).
\]
右侧的同态又零化 \(J^2e_i\)，因而得到内在的识别
\[
\operatorname{Hom}_B(Je_i,S_j)
\cong\operatorname{Hom}_{B/J}(Je_i/J^2e_i,S_j).
\]
这里源模被 \(J\) 零化，是 \(k^r\) 模；映到 \(S_j\) 的同态只保留第 \(j\) 个坐标 \(e_j(J/J^2)e_i\)。\(S_j\) 的指定向量 \(e_j+Je_j\) 非零并生成这个一维空间，按命题中的识别，所列同态恰对应这个坐标上的任意 \(k\) 线性泛函。若只给定与 \(S_j\) 同构的抽象简单模，则写成 \(k\) 对偶需要另选该一维模与 \(k\) 的识别；箭数所用的维数不依赖这一识别。
\end{proof}
这个箭图称为 \(B\) 的\term[Ext-jiantu]{Ext 箭图}[Ext quiver]。箭的具体基依赖选择，但每一对顶点间的箭数由扩张空间决定。注意顺序是 \(\operatorname{Ext}^1(S_i,S_j)\) 对应 \(i\to j\)：其扩张以 \(S_j\) 为子模、\(S_i\) 为商模。

\ProseParagraphBreak

取任意域 \(k\) 上的 \(B=T_2(k)\)，令 \(e_1=E_{11}\)、\(e_2=E_{22}\)。此时 \(J=kE_{12}\)、\(J^2=0\)，而
\[
e_1(J/J^2)e_2=kE_{12},\qquad e_2(J/J^2)e_1=0.
\]
所以唯一的箭为 \(2\to1\)，由\cref{b5:prop:ext-quiver-direction}有
\[
\operatorname{Ext}_B^1(S_2,S_1)\cong k,
\qquad \operatorname{Ext}_B^1(S_1,S_2)=0.
\]
左理想 \(P_2=Be_2=kE_{22}\oplus kE_{12}\) 的根为 \(kE_{12}\cong S_1\)，顶为 \(S_2\)，从而给出
\[
0\longrightarrow S_1\longrightarrow P_2
\longrightarrow S_2\longrightarrow0.
\]
在单箭 \(2\to1\) 的表示中，\(E_{12}E_{22}=E_{12}\) 使中项成为 \(k\xrightarrow{1}k\)，子模与商模上的映射均取恒等。由\cref{b5:exm:arrow-extension-parameter}，这正是非零参数的非分裂扩张，也就是\cref{b5:exm:triangular-projective-covers}中计算的短正合列。箭的终点记录子模，起点记录商模，反向没有对应的非零扩张。

\ProseParagraphBreak

对偶来自扩张参数的取值。设 \(i\ne j\)，并设 \(e_j(J/J^2)e_i\) 以 \(\bar a,\bar b\) 为基，选定提升 \(a,b\in e_jJe_i\)。一个泛函 \(f\) 按前一命题对应于模同态 \(\widetilde f:Je_i\to S_j\)。\Cref{b5:prop:projective-presentation-extensions}给出的推出模是
\[
(S_j\oplus Be_i)/\{(\widetilde f(y),-y):y\in Je_i\}.
\]
被取商的子模使 \((\widetilde f(y),0)\) 与 \((0,y)\) 相差为零，因此有
\[
[(0,xe_i)]=[(\widetilde f(xe_i),0)]\qquad(x\in J).
\]
若 \(u\) 为子模 \(S_j\) 的指定生成元，\(v=[(0,e_i)]\) 为商模 \(S_i\) 的指定生成元的提升，就得到
\[
\begin{aligned}
av&=[(0,ae_i)]=[(\widetilde f(ae_i),0)]=f(\bar a)u,\\
bv&=[(0,be_i)]=[(\widetilde f(be_i),0)]=f(\bar b)u.
\end{aligned}
\]
这个限制映射零化 \(J^2e_i\)，所以只需记录箭的商类。箭基给出两个一次根方向，扩张类则给出在这两个方向上的数值 \(f(\bar a),f(\bar b)\)，这就是线性对偶出现的原因。固定两端时，这些数值是扩张参数；若还允许对子模和商模换基，便进一步改变参数，不能把扩张类与中项同构类混同。

\subsection{容许关系、路径代数与有限维性}
\label{b5:subsec:0905:03}

从根的第一层恢复整个代数，可以先看域 \(k\) 上的上三角矩阵代数 \(T_3(k)\)。它的根及根幂为
\[
J=kE_{12}\oplus kE_{23}\oplus kE_{13},
\qquad J^2=kE_{13},\qquad J^3=0.
\]
在 \(J/J^2\) 中只需选 \(E_{12}+J^2,E_{23}+J^2\) 两个方向，提升的线性生成空间为 \(L=kE_{12}\oplus kE_{23}\)。以 \(L^2\) 表示 \(L\) 中元素两两乘积的线性生成空间，则
\[
E_{13}=E_{12}E_{23},\qquad J^2=L^2,
\qquad J=L+L^2,
\]
较深层的 \(E_{13}\) 已由两个箭像相乘产生，无须再独立选择。一般的证明也逐层把根元素的余项写成更多箭像的乘积，根幂零性保证这个过程有限步结束。

\begin{theorem}\label{b5:thm:basic-algebra-bound-quiver}
设 \(k\) 为域，\(B\) 为非零有限维含幺分裂基本 \(k\) 代数，\(J=J(B)\)，\(B/J\cong k^r\)。取一组完整正交幂等元 \(e_1,\ldots,e_r\)，其模 \(J\) 的像为各标准坐标幂等元。以 \(1,\ldots,r\) 为顶点，对每个 \(e_j(J/J^2)e_i\) 的一组 \(k\) 基中的元素画一支箭 \(i\to j\)，构成有限箭图 \(Q\)。则存在 \(kQ\) 的容许关系理想 \(I\)，使
\[
B\cong kQ/I.
\]
特别地，每个代数闭域上的非零有限维代数都与某个这种商代数 Morita 等价。
\end{theorem}
\begin{proof}
记 \(J=J(B)\)。将顶点送到选定的 \(e_i\)，将每支箭 \(i\to j\) 送到其基元的某个提升 \(x_a\in e_jJe_i\)。顶点关系成立，由路径代数的构造得到同态 \(\phi:kQ\to B\)。令 \(C\) 为这些像生成的子代数，\(L\) 为箭像的线性生成空间。基的选择给出 \(J=L+J^2\)，而顶点像生成 \(B/J\)，故 \(B=C+J\)。
对每个正整数 \(m\)，用 \(L^m\) 表示全部 \(m\) 个 \(L\) 中元素的乘积所线性生成的空间；它是子空间的乘积记号，不要求 \(L\) 是理想。展开 \(m\) 个 \(J=L+J^2\) 的乘积，全部因子来自 \(L\) 的项属于 \(L^m\)，其余项至少有一个 \(J^2\) 因子，故属于 \(J^{m+1}\)。再由 \(L^m,J^{m+1}\subseteq J^m\)，得到
\[
J^m=L^m+J^{m+1}.
\]
取 \(N\ge2\) 使 \(J^N=0\)。任意 \(x\in J\) 先写为 \(x=\ell_1+x_2\)，其中 \(\ell_1\in L\)、\(x_2\in J^2\)；再将 \(x_2\) 写为 \(\ell_2+x_3\)，其中 \(\ell_2\in L^2\)、\(x_3\in J^3\)。继续应用所列等式，最终得到
\[
x=\ell_1+\cdots+\ell_{N-1},\qquad \ell_m\in L^m,
\]
因为最后的余项属于 \(J^N=0\)。所以 \(J=L+\cdots+L^{N-1}\subseteq C\)，进而 \(B=C+J=C\)，\(\phi\) 满射。
令 \(I=\ker\phi\)，并令 \(\mathfrak a\) 为箭理想。因每支箭映入 \(J\)，有 \(\mathfrak a^N\subseteq I\)。为证 \(I\subseteq\mathfrak a^2\)，把核中一个元素拆成长度零项、长度一项及较长路径。其像模 \(J\) 为零，迫使全部顶点系数为零；再模 \(J^2\)，箭像的基独立性迫使全部长度一系数为零。故它只含长度至少二的路径。所求理想容许，第一同构定理给出商表示。最后结合\cref{b5:prop:basic-morita-reduction}，得到一般代数闭基域情形。
\end{proof}
关系理想通常不唯一，因为箭的提升与基可以改变。定理给出的是一种适合模范畴的展示，并不宣称所有代数都有唯一的一组路径关系。这里的生成性也用到了根的幂零性；若只知道根的第一层而没有这种有限性，就不能用同样的有限步推导结束。

\subsection{三角矩阵、双数与带关系代数}
\label{b5:subsec:0905:04}

前面的 \(T_3(k)\) 计算给出箭图 \(3\to2\to1\)，两支箭分别对应 \(E_{23},E_{12}\)。它的六条路径对应六个上三角矩阵单位，长度二路径的像为 \(E_{12}E_{23}=E_{13}\ne0\)，因此路径代数已经同构于 \(T_3(k)\)，不必再加入关系。这也复核了左模约定下的箭向。
\ProseParagraphBreak
双数代数 \(k[t]/(t^2)\) 则只有一个顶点和一支圈，关系为圈的平方为零。它的简单模 \(S=k\) 满足
\[
\operatorname{Ext}^1(S,S)\cong k,
\]
因为 \(J/J^2=kt\)。同样的单圈若改加 \(t^3=0\)，Ext 箭图仍相同，代数维数却从二变成三，最长不可分解的幂零链也从二变成三。箭图本身没有记录全部乘法，关系不可省去。
\begin{example}\label{b5:exm:zero-composition-bound-quiver}
设 \(k\) 为域，\(Q=(1\xrightarrow{a}2\xrightarrow{b}3)\)，\(B=kQ/(ba)\)，并沿用同一记号表示路径的商类。其基为 \(e_1,e_2,e_3,a,b\)，根为 \(J=ka\oplus kb\)，满足 \(J^2=0\)，Ext 箭图仍是原来的三顶点直线。然而维数向量 \((1,1,1)\) 的 \(B\) 表示均可分解，而无关系的 \(kQ\) 在这个维数向量上有不可分解表示 \(I[1,3]\)。令 \(P_i=Be_i\)、\(S_1=P_1/JP_1\)，则下列左 \(B\)-模序列是长度二的投射消解：
\[
0\longrightarrow P_3\xrightarrow{\;q\mapsto qb\;}P_2
\xrightarrow{\;p\mapsto pa\;}P_1
\longrightarrow S_1\longrightarrow0.
\]
\end{example}
\begin{proof}
理想 \(ka\oplus kb\) 的平方为零，商为 \(k^3\)，所以它就是 \(J(B)\)。两支箭在 \(J/J^2\) 中仍独立，按\cref{b5:prop:ext-quiver-direction}得到原箭图。维数向量 \((1,1,1)\) 的表示中，两支箭为标量 \(\alpha,\beta\)，关系要求 \(\beta\alpha=0\)。至少一支箭为零，表示便沿该处拆成两个非零直和项；无关系时，两箭均取恒等的 \(I[1,3]\) 不可分解。

三个左理想的路径基为
\[
P_1=ke_1\oplus ka,\qquad
P_2=ke_2\oplus kb,\qquad P_3=ke_3.
\]
它们都是左正则模的直和因子，故投射。右乘 \(a\) 与右乘 \(b\) 都给出左 \(B\) 线性同态，路径乘法具体为
\[
e_2a=a,\qquad ba=0,\qquad e_3b=b.
\]
所以 \(P_2\to P_1\) 的像为 \(ka=JP_1\)，正是自然商 \(P_1\to S_1\) 的核；其核为 \(kb\)，又正是单射 \(P_3\to P_2\) 的像。这逐项证明了所列序列正合。

为比较无关系情形，记 \(P_i^0=kQe_i\)，把同一个顶点简单模 \(S_1\) 看成左 \(kQ\)-模。此时 \(P_1^0\) 以 \(e_1,a,ba\) 为基，\(P_2^0\) 以 \(e_2,b\) 为基。右乘 \(a\) 将后一组基送到 \(a,ba\)，两条路径独立，因而给出长度一投射消解
\[
0\longrightarrow P_2^0\xrightarrow{\;p\mapsto pa\;}P_1^0
\longrightarrow S_1\longrightarrow0.
\]
取商使 \(ba\) 从非零路径变为零，原先单射的右乘映射便新增核 \(kb\)。所列长度二消解中的 \(P_3\) 正是用来覆盖这个由关系产生的核。
\end{proof}

\subsection{ADE 二次型与反射简介}
\label{b5:subsec:0905:05}

对有限无环箭图，\cref{b5:thm:quiver-hom-ext-euler}已经给出
\[
q_Q(\mathbf d)=\sum_i d_i^2-\sum_{a:i\to j}d_i d_j
=\dim\operatorname{End}_Q(V)-\dim\operatorname{Ext}_Q^1(V,V),
\quad \mathbf d=\mathbf d(V).
\]
右边来自表示的线性方程，左边却只用图和整数向量。在线性定向 \(A_n\) 中，一个区间向量含有 \(b-a+1\) 个一，区间内有 \(b-a\) 支箭，所以其值为一。我们已经用选基法证明了这些区间给出全部不可分解表示；这个数值因而有实际分类意义。
\ProseParagraphBreak
其他无向图是否也有类似现象，需要继续研究 \(q_Q\) 的正定性以及改变箭向的办法。\chapref{b5:ch:10}将定义整数格上的反射，证明所需的 ADE 图判别，再构造反射函子。对于带关系代数，\cref{b5:exm:zero-composition-bound-quiver}已经说明，仅仅保留同一个图上的二次式还不足以恢复不可分解表示。
